\section{Our Approach}
\label{sec:impl}

All five implementations grow the same exact CART or C4.5 tree.
They differ in how the sorted order of a node's samples is obtained, in how
many candidate cuts are scored exactly, and in where and in what order the
nodes are grown.
\ImplGpu{} and \ImplHybrid{} are two configurations of one CUDA code path:
the hand-off threshold is one sample, so every node stays on the GPU, or
small subtrees go to the CPU.

\subsection{Common Design}
\label{sec:impl-common}

The training set is stored column-major as 32-bit floats, and labels are
mapped to dense integers at load time.
Every implementation scores a cut with the same double-precision routine,
compiled for the CPU and for the GPU.
Gains within \(10^{-12}\) are ties: the earlier cut wins, then the lower
attribute index.
Filters discard only cuts that cannot replace the incumbent; the cut that is
kept is scored by the shared routine.

\ImplSerialBase{} works on lists of sample indices.
The other four store each attribute as 8-byte entries, the value together
with the sample and the class.
A node occupies the same index range in every column.
A stable partition of that range leaves both children sorted, so each
attribute is ordered only once: the presort of SLIQ and
SPRINT~\cite{mehta1996sliq,shafer1996sprint}, with the class kept in the
entry.
The tree is a flat preorder array.
Concurrent builders write nodes to a shared store and number them afterwards.
Our tests require identical trees from all five implementations; the checks
against scikit-learn and Quinlan's program are those of
Section~\ref{sec:intro}.
Pruning, when used, runs afterwards on the CPU.
The timings in Section~\ref{sec:design} are depth-limited trees left
unpruned.

\subsection{\ImplSerialBase}
\label{sec:impl-simple}

\ImplSerialBase{} is the textbook procedure on one thread, and the reference
for the others.
At each node it sorts the node's sample indices by every attribute, scores
every admissible cut, and recurses on two new lists.
A node of \(n\) samples and \(F\) attributes costs \(O(Fn\log n)\), and the
sort reads values through the indices.

\subsection{\ImplSerialOpt}
\label{sec:impl-serial}

\ImplSerialOpt{} is the builder of Section~\ref{sec:impl-parallel} with the
thread pool off, so the presort and the growth use one thread.
A stable radix sort orders each attribute once.
The left child is the first \(n_L\) entries of the winning column; the other
columns are partitioned stably to match, without comparing values with the
threshold again.
An attribute constant on the node is dropped for the whole subtree.
The sweep keeps a cheap estimate of the gain, and a convex bound for two
classes.
Only cuts that can still win are scored by the shared routine, in order.
If both children are already leaves, the columns are not moved.

The columns occupy \(F+1\) slots of \(8N\) bytes, allocated once.
Nodes of at least \(2^{15}\) samples move out of place into the spare slot,
through a one-bit map of the selected rows.
Smaller nodes are partitioned in place through a reused buffer.
Large allocations request transparent huge pages, and an explicit stack
replaces recursion.

\subsection{\ImplParallel}
\label{sec:impl-parallel}

\ImplParallel{} runs this builder on a thread pool and, as in the mixed
strategy of Aldinucci et al.~\cite{aldinucci2014fastflow}, chooses the
parallelism by node size.
On a node of at least 32{,}768 samples, attribute scans run together and the
partition uses one column per thread; below that, waking the pool costs more
than the scan.
After a node of at least 4{,}096 samples is split, the left child becomes a
task and the thread continues with the right child.
The pool joins once, at the end.
Up to \(\min(F-1, p)\) spare slots, for \(p\) threads, let a large node move
several columns at once, and the presort uses the same pool.

The CPU builder does not pin workers or bind memory.
The thread count follows the process affinity mask, so \texttt{numactl} or
\texttt{taskset} can confine a run, and each page stays on the NUMA node of
the thread that first writes it.

\subsection{\ImplGpu}
\label{sec:impl-gpu}

\ImplGpu{} grows the tree breadth-first, because a small node cannot fill the
device.
Every node of a level is handled together.
Work is cut into tiles of 2{,}048 entries of one node's range in one column,
one thread block per tile.
The device holds the columns in \(F+1\) slots and radix-sorts each attribute,
stably, as soon as it has been uploaded.
The best cut of every attribute returns to the host, which applies the same
rule as the CPU and closes children that are already leaves: the only
decision round trip of the level.
A single-pass stable partition writes the children into a free slot and
counts the next level.
On the A100 and the H100 of Section~\ref{sec:design}, every admissible cut is
scored in double precision.
Where double precision is much slower, a single-precision pass first drops
cuts that cannot win.

The training matrix stays pageable and is copied through a few page-locked
staging buffers while earlier chunks are in flight.
When the process may run outside the GPU's NUMA node, those buffers are
allocated on that node and the copy threads are pinned to its CPUs.
This is the only explicit NUMA placement in the source; \ImplHybrid{} pins
its copy threads in the same way.
The training columns and the worker pool are not bound to that node.

\subsection{\ImplHybrid}
\label{sec:impl-hybrid}

\ImplHybrid{} assigns nodes of this one tree by size
(Fig.~\ref{fig:approach-hybrid}).
A GPU level costs about as much for many small nodes as for a few large
ones, while the CPU grows a small subtree as an independent task.
A child that still needs a split and has fewer than \(m\) samples is handed
to the builder of Section~\ref{sec:impl-parallel}; the GPU continues.
By default \(m\) is \(N/8192\), clamped to \([1024, 2048]\), for CART and
\(N/4096\), clamped to \([1024, 4096]\), for C4.5, bounds chosen empirically
and overridable.

\begin{figure}[t]
\centering
\includegraphics[width=\textwidth]{figs/approach-hybrid.pdf}
\caption{One level of \ImplHybrid{}.
The GPU sweeps and partitions the large nodes together; the host chooses
each split.
Children with fewer than \(m\) samples are gathered into a device ring,
copied through pinned buffers into a pageable host arena, and grown by the
CPU pool while the GPU continues with the next level.
Both paths write one node store.
\ImplGpu{} is the same loop with \(m = 1\).}
\label{fig:approach-hybrid}
\end{figure}

The gather keeps only attributes that can still split the child, and writes
them into a device ring.
Copy streams move the data through two pinned buffers into a pageable host
arena.
A host callback starts the CPU tasks when the copy arrives, so the GPU does
not wait for them; it waits on the ring only when it is a whole ring ahead
of the link.
CPU threads touch the arena in the background.
The arena is not allocated on the GPU's NUMA node.
Both devices append to the store of Section~\ref{sec:impl-common}.
